\documentclass[11pt,letterpaper]{article}
\usepackage[utf8]{inputenc}
\usepackage{geometry}
\usepackage{setspace}

\geometry{letterpaper, margin=1in}
\doublespacing

\setlength{\parindent}{0.5in}
\setlength{\parskip}{0pt}

\begin{document}

\thispagestyle{empty}

\begin{center}
    \textbf{\large ABSTRACT OF THE DISCLOSURE}
\end{center}

\vspace{0.6cm}

A computing system that reduces or eliminates the need for a global, high-frequency clock signal by using an internal, massively parallel optical data fabric for communication. The system comprises multiple processing components (e.g., CPU cores, accelerators) interconnected by optical buses. Instead of operating on a strict, system-wide clock cycle, components operate asynchronously, triggered by the arrival of data frames over the optical fabric. The ultra-high bandwidth and low latency of the optical connections ensure that data arrives faster than a typical clock cycle, making a data-flow or event-driven architecture more efficient and eliminating the power consumption and design complexity associated with distributing a high-frequency clock signal across a chip or system.

\end{document}
